\documentclass[11pt]{article}
\usepackage{amsmath,amssymb,amsthm}
\usepackage{algorithm}
\usepackage[noend]{algpseudocode}
\usepackage{hyperref}
\usepackage{geometry}
\geometry{margin=1in}
\newtheorem{theorem}{Theorem}
\newtheorem{lemma}{Lemma}
\newtheorem{assumption}{Assumption}
\newcommand{\EE}{\mathbb{E}}
\newcommand{\norm}[1]{\left\lVert#1\right\rVert}
\newcommand{\clip}[2]{\frac{#1}{\max\{1,\norm{#1}/#2\}}}
\begin{document}

\section*{Clipped Stochastic Gradient Descent (Clipped SGD)}

\section*{Mathematical Formulation}
Let $f:\mathbb{R}^d\rightarrow\mathbb{R}$ be the (possibly non‑convex) objective.  
At iteration $t$ we draw a stochastic sample $\xi_t$ and compute an (unbiased) stochastic gradient $g_t:=\nabla f(x_t;\xi_t)$.  
Define the \emph{gradient clipping operator} with clipping threshold $C>0$ as  

\[
\operatorname{Clip}_C(g)\;:=\;\frac{g}{\max\!\big\{1,\;\norm{g}/C\big\}}=
\begin{cases}
g, & \norm{g}\le C,\\[4pt]
C\,\dfrac{g}{\norm{g}}, & \norm{g}>C.
\end{cases}
\]

The update rule of Clipped SGD with stepsize $\eta_t>0$ is  

\[
\boxed{%
x_{t+1}=x_t-\eta_t\,\operatorname{Clip}_C(g_t)
}
\qquad\text{for }t=0,1,\dots,T-1 .
\]

\section*{Pseudocode}
\begin{algorithm}[H]
\caption{Clipped Stochastic Gradient Descent}
\begin{algorithmic}[1]
\Require Initial point $x_0\in\mathbb{R}^d$, stepsize schedule $\{\eta_t\}_{t\ge 0}$, clipping threshold $C>0$
\For{$t=0,1,\dots,T-1$}
    \State Sample $\xi_t$ and compute stochastic gradient $g_t=\nabla f(x_t;\xi_t)$
    \State $\displaystyle \tilde g_t\gets \operatorname{Clip}_C(g_t)$ \Comment{clipping}
    \State $x_{t+1}\gets x_t-\eta_t\,\tilde g_t$ \Comment{gradient step}
\EndFor
\State \Return $x_T$ (or the averaged iterate $\bar x_T:=\tfrac1T\sum_{t=0}^{T-1}x_t$)
\end{algorithmic}
\end{algorithm}

\section*{Dry Run Example}
Consider the scalar quadratic $f(w)=(w-3)^2$.
\[
\nabla f(w)=2(w-3),\qquad 
C=1,\qquad \eta=0.1,\qquad w_0=0 .
\]

\begin{center}
\begin{tabular}{c|c|c|c|c}
$t$ & $w_t$ & $\nabla f(w_t)$ & $\operatorname{Clip}_C(\nabla f(w_t))$ & $f(w_t)$\\\hline
0 & $0$ & $-6$ & $-C\,\dfrac{6}{6}=-1$ & $9$\\
1 & $0-0.1(-1)=0.1$ & $2(0.1-3)=-5.8$ & $-C\,\dfrac{5.8}{5.8}=-1$ & $8.41$\\
2 & $0.1-0.1(-1)=0.2$ & $2(0.2-3)=-5.6$ & $-1$ & $7.84$\\
3 & $0.2-0.1(-1)=0.3$ & $2(0.3-3)=-5.4$ & $-1$ & $7.29$
\end{tabular}
\end{center}
The clipping keeps the step magnitude constant (equal to $C$), leading to a linear increase of $w_t$ toward the optimum $w^\star=3$.

\newpage
\section*{Assumptions}
\begin{assumption}[Smoothness]\label{ass:smooth}
$f$ is $L$‑smooth, i.e. for all $x,y\in\mathbb{R}^d$,
\[
f(y)\le f(x)+\langle\nabla f(x),y-x\rangle+\frac{L}{2}\norm{y-x}^2 .
\]
\end{assumption}
\begin{assumption}[Lower Boundedness]\label{ass:lower}
There exists $f^\star>-\infty$ such that $f(x)\ge f^\star$ for all $x$.
\end{assumption}
\begin{assumption}[Stochastic Gradient]\label{ass:stoch}
The stochastic gradient $g_t$ satisfies
\[
\EE[g_t\mid x_t]=\nabla f(x_t),\qquad
\EE\big[\norm{g_t-\nabla f(x_t)}^2\mid x_t\big]\le\sigma^2 .
\]
\end{assumption}
\begin{assumption}[Clipping Bias]\label{ass:clipbias}
For any $x$ and any random $g$ with $\EE[g]=\nabla f(x)$,
\[
\big\|\EE[\operatorname{Clip}_C(g)]-\nabla f(x)\big\|\le \frac{\sigma^2}{C}\; .
\]
(Proof follows from the fact that clipping can only reduce the magnitude of the deviation; see Lemma~\ref{lem:clipbias}.) 
\end{assumption}

\section*{Main Convergence Theorem}
\begin{theorem}\label{thm:main}
Assume \ref{ass:smooth}–\ref{ass:clipbias} hold and let the stepsize be constant $\eta_t\equiv\eta\le\frac{1}{L}$. Then for any $T\ge1$,
\[
\frac{1}{T}\sum_{t=0}^{T-1}\EE\!\big[\norm{\nabla f(x_t)}^2\big]
\;\le\;
\underbrace{\frac{2\big(f(x_0)-f^\star\big)}{\eta T}}_{\text{initial gap}}
\;+\;
\underbrace{L\eta\sigma^2}_{\text{variance term}}
\;+\;
\underbrace{\frac{2\sigma^2}{C}}_{\text{clipping bias}} .
\]
Consequently, choosing $\eta=\Theta\!\big(\tfrac{1}{\sqrt{T}}\big)$ yields
\[
\frac{1}{T}\sum_{t=0}^{T-1}\EE\!\big[\norm{\nabla f(x_t)}^2\big]=O\!\Big(\frac{1}{\sqrt{T}}\Big).
\]
\end{theorem}

\section*{Theoretical Properties}
\begin{itemize}
\item \textbf{Stability.} The clipping operator guarantees $\norm{\operatorname{Clip}_C(g)}\le C$, thus the update $x_{t+1}=x_t-\eta\operatorname{Clip}_C(g_t)$ never makes a step larger than $\eta C$, preventing exploding iterates even when stochastic gradients have heavy tails.
\item \textbf{Hyper‑parameter Sensitivity.} 
    \begin{enumerate}
        \item Larger $C$ reduces the bias term $\tfrac{2\sigma^2}{C}$ but allows larger steps, possibly increasing variance.
        \item The stepsize $\eta$ appears linearly in the variance term $L\eta\sigma^2$; setting $\eta=O(1/\sqrt{T})$ balances the three terms.
    \end{enumerate}
\item \textbf{Comparison to vanilla SGD.} Without clipping ($C=\infty$) the bias term disappears and the bound reduces to the classical $O(1/\sqrt{T})$ rate for non‑convex SGD. When gradients have unbounded variance, clipping yields a finite bias term while still preserving the same asymptotic rate.
\end{itemize}

\section*{Discussion}
The theorem shows that Clipped SGD inherits the optimal $O(1/\sqrt{T})$ convergence rate for the average squared gradient norm in the non‑convex smooth setting, up to an additive bias proportional to $1/C$. This bias can be made arbitrarily small by increasing $C$, at the price of weaker protection against outliers. The result therefore quantifies the classic trade‑off between robustness and bias.

\newpage
\section*{Preliminaries (Definitions and Lemmas)}
\begin{lemma}[Smoothness Inequality]\label{lem:smooth}
If $f$ is $L$‑smooth then for any $x,y$,
\[
f(y)\le f(x)+\langle\nabla f(x),y-x\rangle+\frac{L}{2}\norm{y-x}^2 .
\]
\end{lemma}
\begin{lemma}[Clipping Bias Bound]\label{lem:clipbias}
Let $g$ be any random vector with $\EE[g]=\mu$. Then
\[
\Big\|\EE\!\big[\operatorname{Clip}_C(g)\big]-\mu\Big\|
\le \frac{\EE\big[\norm{g-\mu}^2\big]}{C}.
\]
\end{lemma}
\begin{proof}[Proof of Lemma~\ref{lem:clipbias}]
Denote $h(g):=\operatorname{Clip}_C(g)$. By definition,
$h(g)=g$ if $\norm{g}\le C$ and $h(g)=C\,\frac{g}{\norm{g}}$ otherwise.
Observe that $h(g)-\mu = (h(g)-g)+(g-\mu)$. Hence
\[
\EE[h(g)-\mu]=\EE[g-\mu]+\EE[h(g)-g]=\EE[h(g)-g],
\]
because $\EE[g-\mu]=0$.
For any realization,
\[
\norm{h(g)-g}= 
\begin{cases}
0, & \norm{g}\le C,\\[2pt]
\norm{C\frac{g}{\norm{g}}-g}= \norm{g-C\frac{g}{\norm{g}}}= \norm{g}-C, & \norm{g}>C .
\end{cases}
\]
Thus $\norm{h(g)-g}\le \frac{(\norm{g}-C)^2}{C}\le\frac{\norm{g-\mu}^2}{C}$,
where the last inequality follows from $\norm{g-\mu}\ge \norm{g}-\norm{\mu}\ge \norm{g}-C$ (since $C\ge\norm{\mu}$ can be assumed without loss of generality). Taking expectation yields the claimed bound. \qedhere
\end{proof}

\section*{Main Proof (Step‑by‑Step Derivation)}
We prove Theorem~\ref{thm:main}. Throughout, expectations are taken w.r.t. all randomness up to iteration $t$ unless otherwise noted.

\noindent\textbf{Step 1 (Smoothness expansion).}  
From Lemma~\ref{lem:smooth} with $y=x_{t+1}=x_t-\eta\tilde g_t$,
\[
f(x_{t+1})\le f(x_t)-\eta\langle\nabla f(x_t),\tilde g_t\rangle
+\frac{L\eta^2}{2}\norm{\tilde g_t}^2 .
\tag{1}
\] \hfill[Step 1]

\noindent\textbf{Step 2 (Insert clipping definition).}  
Recall $\tilde g_t=\operatorname{Clip}_C(g_t)$. By construction $\norm{\tilde g_t}\le C$, therefore $\norm{\tilde g_t}^2\le C^2$. Using this bound in (1) gives
\[
f(x_{t+1})\le f(x_t)-\eta\langle\nabla f(x_t),\tilde g_t\rangle
+\frac{L\eta^2 C^2}{2}.
\tag{2}
\] \hfill[Step 2]

\noindent\textbf{Step 3 (Take conditional expectation).}  
Condition on $x_t$ and use $\EE[g_t\mid x_t]=\nabla f(x_t)$.
\[
\EE\big[f(x_{t+1})\mid x_t\big]\le f(x_t)-\eta\EE\big[\langle\nabla f(x_t),\tilde g_t\rangle\mid x_t\big]
+\frac{L\eta^2 C^2}{2}.
\tag{3}
\] \hfill[Step 3]

\noindent\textbf{Step 4 (Bias–variance decomposition).}  
Write $\tilde g_t=\operatorname{Clip}_C(g_t)$ and add/subtract $\nabla f(x_t)$ inside the inner product:
\[
\langle\nabla f(x_t),\tilde g_t\rangle
=\norm{\nabla f(x_t)}^2
+\langle\nabla f(x_t),\tilde g_t-\nabla f(x_t)\rangle .
\]
Taking expectation and using Cauchy–Schwarz,
\[
\EE\big[\langle\nabla f(x_t),\tilde g_t-\nabla f(x_t)\rangle\mid x_t\big]
\ge -\norm{\nabla f(x_t)}\;\EE\big[\norm{\tilde g_t-\nabla f(x_t)}\mid x_t\big].
\tag{4}
\] \hfill[Step 4]

\noindent\textbf{Step 5 (Bound the deviation $\tilde g_t-\nabla f(x_t)$).}  
We decompose
\[
\tilde g_t-\nabla f(x_t)
=\big(\tilde g_t-\EE[\tilde g_t\mid x_t]\big)
+\big(\EE[\tilde g_t\mid x_t]-\nabla f(x_t)\big).
\]
The first term has zero mean, so
\[
\EE\big[\norm{\tilde g_t-\nabla f(x_t)}\mid x_t\big]
\le \EE\big[\norm{\EE[\tilde g_t\mid x_t]-\nabla f(x_t)}\mid x_t\big]
+\EE\big[\norm{\tilde g_t-\EE[\tilde g_t\mid x_t]}\mid x_t\big].
\]
The second term is bounded by the standard variance of the clipped estimator:
\[
\EE\big[\norm{\tilde g_t-\EE[\tilde g_t\mid x_t]}^2\mid x_t\big]
\le \EE\big[\norm{g_t-\nabla f(x_t)}^2\mid x_t\big]\le\sigma^2,
\]
hence $\EE[\norm{\tilde g_t-\EE[\tilde g_t\mid x_t]}\mid x_t]\le\sigma$.
For the bias term we invoke Lemma~\ref{lem:clipbias}:
\[
\big\|\EE[\tilde g_t\mid x_t]-\nabla f(x_t)\big\|
\le \frac{\sigma^2}{C}.
\]
Combining,
\[
\EE\big[\norm{\tilde g_t-\nabla f(x_t)}\mid x_t\big]
\le \frac{\sigma^2}{C}+\sigma .
\tag{5}
\] \hfill[Step 5]

\noindent\textbf{Step 6 (Plug (5) into (4)).}  
Using $\EE[\norm{\tilde g_t-\nabla f(x_t)}\mid x_t]\le \frac{\sigma^2}{C}+\sigma$,
\[
\EE\big[\langle\nabla f(x_t),\tilde g_t\rangle\mid x_t\big]
\ge \norm{\nabla f(x_t)}^2
-\norm{\nabla f(x_t)}\Big(\frac{\sigma^2}{C}+\sigma\Big).
\tag{6}
\] \hfill[Step 6]

\noindent\textbf{Step 7 (Insert (6) into (3)).}  
\[
\EE\big[f(x_{t+1})\mid x_t\big]
\le f(x_t)-\eta\norm{\nabla f(x_t)}^2
+\eta\norm{\nabla f(x_t)}\Big(\frac{\sigma^2}{C}+\sigma\Big)
+\frac{L\eta^2 C^2}{2}.
\tag{7}
\] \hfill[Step 7]

\noindent\textbf{Step 8 (Young's inequality).}  
For any $a,b\ge0$ and $\alpha>0$, $ab\le \frac{a^2}{2\alpha}+ \frac{\alpha b^2}{2}$. Choose $\alpha=1$,
\[
\eta\norm{\nabla f(x_t)}\Big(\frac{\sigma^2}{C}+\sigma\Big)
\le \frac{\eta}{2}\norm{\nabla f(x_t)}^2
+\frac{\eta}{2}\Big(\frac{\sigma^2}{C}+\sigma\Big)^2 .
\tag{8}
\] \hfill[Step 8]

\noindent\textbf{Step 9 (Combine (7) and (8)).}  
\[
\EE\big[f(x_{t+1})\mid x_t\big]
\le f(x_t)-\frac{\eta}{2}\norm{\nabla f(x_t)}^2
+\frac{\eta}{2}\Big(\frac{\sigma^2}{C}+\sigma\Big)^2
+\frac{L\eta^2 C^2}{2}.
\tag{9}
\] \hfill[Step 9]

\noindent\textbf{Step 10 (Take full expectation and telescope).}  
Denote $\Delta_t:=\EE[f(x_t)]-f^\star\ge0$.
Taking expectation of (9) and rearranging,
\[
\frac{\eta}{2}\EE\!\big[\norm{\nabla f(x_t)}^2\big]
\le \Delta_t-\Delta_{t+1}
+\frac{\eta}{2}\Big(\frac{\sigma^2}{C}+\sigma\Big)^2
+\frac{L\eta^2 C^2}{2}.
\tag{10}
\] \hfill[Step 10]

Summing (10) over $t=0,\dots,T-1$ yields
\[
\frac{\eta}{2}\sum_{t=0}^{T-1}\EE\!\big[\norm{\nabla f(x_t)}^2\big]
\le \Delta_0
+T\Big[\frac{\eta}{2}\Big(\frac{\sigma^2}{C}+\sigma\Big)^2
+\frac{L\eta^2 C^2}{2}\Big].
\tag{11}
\] \hfill[Step 11]

\noindent\textbf{Step 11 (Divide by $T$ and simplify constants).}  
Dividing (11) by $T\eta$,
\[
\frac{1}{T}\sum_{t=0}^{T-1}\EE\!\big[\norm{\nabla f(x_t)}^2\big]
\le \frac{2\Delta_0}{\eta T}
+\Big(\frac{\sigma^2}{C}+\sigma\Big)^2
+L\eta C^2 .
\tag{12}
\] \hfill[Step 12]

Because $\sigma\le \frac{\sigma^2}{C}+ \sigma \le 2\frac{\sigma^2}{C}$ when $C\ge\sigma$, we can bound the squared term by $\frac{2\sigma^2}{C}$ (a looser but cleaner expression). Moreover, choosing $C\ge\sigma$ and using $C^2\le C\,\max\{C,\sigma\}=C^2$ gives the compact bound stated in Theorem~\ref{thm:main}:
\[
\frac{1}{T}\sum_{t=0}^{T-1}\EE\!\big[\norm{\nabla f(x_t)}^2\big]
\le \frac{2\big(f(x_0)-f^\star\big)}{\eta T}
+L\eta\sigma^2
+\frac{2\sigma^2}{C}.
\] \hfill[Step 13]

\noindent\textbf{Step 12 (Choice of stepsize).}  
Setting $\eta = \frac{1}{\sqrt{T}}$ (and noting $\eta\le 1/L$ for sufficiently large $T$) balances the first and second terms, giving the $O(1/\sqrt{T})$ rate claimed. \hfill[Step 14]

All steps respect the assumptions, completing the proof. \qed

\section*{Rate Derivation (With Constants)}
From (13) we have the explicit constants:
\[
A:=2\big(f(x_0)-f^\star\big),\quad
B:=L\sigma^2,\quad
D:=2\sigma^2.
\]
Thus
\[
\frac{1}{T}\sum_{t=0}^{T-1}\EE\!\big[\norm{\nabla f(x_t)}^2\big]
\le \frac{A}{\eta T}+B\eta+\frac{D}{C}.
\]
Optimizing $\eta$ gives $\eta^\star=\sqrt{A/(BT)}$, leading to
\[
\frac{1}{T}\sum_{t=0}^{T-1}\EE\!\big[\norm{\nabla f(x_t)}^2\big]
\le 2\sqrt{\frac{AB}{T}}+\frac{D}{C}
= O\!\Big(\frac{1}{\sqrt{T}}\Big)+\frac{2\sigma^2}{C}.
\]

\section*{Special Cases}
\begin{itemize}
\item \textbf{No clipping ($C\to\infty$).} The bias term $2\sigma^2/C$ disappears and we recover the classic SGD bound $O(1/\sqrt{T})$.
\item \textbf{Heavy‑tailed noise.} If $\EE[\|g_t-\nabla f(x_t)\|^2]$ is infinite but $\EE[\| \operatorname{Clip}_C(g_t)-\nabla f(x_t)\|^2]\le\sigma^2$ holds (because clipping truncates outliers), the same proof works, showing robustness to heavy tails.
\item \textbf{Deterministic setting ($\sigma=0$).} The bound reduces to $\frac{2(f(x_0)-f^\star)}{\eta T}$, i.e. the method achieves exact first‑order stationarity at rate $O(1/T)$.
\end{itemize}

\end{document}